\documentclass[11pt]{article}
\usepackage[margin=1in]{geometry}
\usepackage{amsmath,amssymb}
\usepackage{natbib}
\usepackage{graphicx}
\usepackage{booktabs}
\usepackage{multirow}
\usepackage{placeins}
\usepackage[font=small,labelfont=bf]{caption}
\usepackage[colorlinks=true,allcolors=blue]{hyperref}
\setlength{\emergencystretch}{3em}

\title{\bfseries Pockmark fields are locally ordered but not hyperuniform:\\
a structure-factor analysis of Mediterranean seabed fluid-escape depressions}
\author{Ruqing Chen\\[2pt]
\small GUT Geoservice Inc., Montr\'eal, Qu\'ebec, Canada\\
\small \texttt{ruqing.chen@gutgeoservice.com}}
\date{\today}

\begin{document}
\maketitle

\begin{abstract}
\noindent
Submarine pockmarks --- seabed depressions left by ascending fluids --- occur in
fields whose spatial arrangement has been described both as fault-controlled
clustering and as regular spacing set by lateral inhibition. We ask a sharper,
previously untested question: are pockmark fields \emph{hyperuniform} --- that
is, do their long-wavelength density fluctuations vanish, as in avian
photoreceptor mosaics and other disordered hyperuniform patterns --- or is any
regularity purely local? Using single fields drawn from a comprehensive open
geodatabase of Mediterranean pockmarks, we combine the Clark--Evans
nearest-neighbour index (local order) with the radially-averaged structure
factor $S(k)$ and the number variance $\sigma^2(R)$ (long-range order), each
compared against complete-spatial-randomness, hard-core (random sequential
adsorption) and inhomogeneous-Poisson nulls. Local order spans the full range:
one field is strongly clustered (nearest-neighbour index $R=0.47$, below its
heterogeneity null), one is essentially random ($R=1.08$), and one is genuinely
regular ($R=1.26$, exceeding even a hard-core gas). Yet \emph{every} field is
long-range super-Poissonian --- $S(k\!\to\!0)>1$ and $\sigma^2(R)$ growing with
window area, not perimeter --- and in each case the long-range behaviour matches
the field's own inhomogeneous-Poisson null. No field is hyperuniform; the single
regular field is an anisotropic lineament, not an isotropic pattern. Pockmark
order is therefore local and steric or structural, and is masked at long range
by density heterogeneity. The result places pockmark fields with other
self-organised geomorphic patterns whose local order does not extend to
hyperuniformity, and shows that nearest-neighbour regularity must not be read as
long-range order without a structure-factor test.
\end{abstract}

\section{Introduction}

Pockmarks are circular to elliptical seabed depressions formed where gas or pore
water escapes through fine-grained sediment, first described on the Scotian
Shelf \citep{king1970} and now known from continental shelves, slopes, fjords
and lakes worldwide \citep{hovland1988}. Within a field, individual pockmarks
number from tens to many thousands, and the geometry of their spatial
arrangement is a long-standing question. Two end-member pictures recur. In the
first, pockmarks track subsurface fluid plumbing --- faults, polygonal-fault
networks, gas chimneys --- and are correspondingly aligned or clustered. In the
second, a field self-organises: each pockmark draws fluid from, and so depletes,
a surrounding area, inhibiting the formation of neighbours within a
characteristic distance, so that the field acquires a regular, quasi-periodic
spacing. This lateral-inhibition picture, and the use of nearest-neighbour and
Ripley's-$K$ statistics to detect the resulting regularity, was set out by
\citet{hammer2009} for a North Sea field.

A distinct and stronger notion of spatial order has meanwhile become central in
condensed-matter physics and is increasingly applied to biological and
geological patterns: \emph{hyperuniformity}. A point pattern is hyperuniform if
its large-scale density fluctuations are anomalously suppressed --- formally, if
the structure factor $S(k)\!\to\!0$ as $k\!\to\!0$, equivalently if the variance
of the number of points in a window of radius $R$ grows more slowly than the
window's area (as its perimeter, $\sigma^2(R)\sim R$, rather than $\sigma^2(R)
\sim R^2$) \citep{torquato2003,torquato2018}. Hyperuniformity is a \emph{global}
property: it cannot be read off local spacing. Disordered hyperuniform patterns
occur in avian photoreceptor mosaics \citep{jiao2014} and a growing list of
biological and physical systems, where they are thought to reflect an optimised
balance of local irregularity and long-range uniformity.

The two notions are easily conflated. A pattern with a tight nearest-neighbour
spacing --- a hard-core liquid, or a field shaped by lateral inhibition --- has a
nearest-neighbour index well above the random value and looks ``regular'', yet
its long-wavelength fluctuations need not be suppressed at all: hard-core
liquids are \emph{not} hyperuniform. Distinguishing the two requires a
long-range diagnostic (the structure factor or number variance), not a local
one. Here we apply this distinction to pockmark fields for the first time. Using
single fields from a comprehensive open geodatabase of Mediterranean pockmarks
\citep{spatola2025}, we ask whether the spacing regularity reported for such
fields amounts to hyperuniformity, or is purely local order superimposed on a
heterogeneous, fluid- and fault-controlled background.

\section{Data and methods}

\paragraph{Data.} We use the geodatabase of \citet{spatola2025}, which compiles
$7{,}516$ georeferenced pockmarks across the Mediterranean and Black Seas from
the published literature, each with up to $31$ attributes, distributed as an
ESRI point shapefile (projected coordinate system WGS84 UTM zone 33N, metres;
licence CC-BY-4.0). Because the database is a multi-source compilation, we do not
analyse it as a whole --- pooling independent fields would trivially drive any
statistic toward the uncorrelated (Poisson) limit. Instead we segment it by the
\texttt{Location} attribute and select single fields that are (i) spatially
contiguous, (ii) sufficiently populated, and (iii) free of the duplicate- and
rounded-coordinate artefacts present in some entries. Three fields, spanning the
range of local order, are analysed in detail: \emph{Offshore Zannone Island}
($N=112$), the \emph{northwestern Calabrian margin} ($N=418$), and the
\emph{Malta Plateau} ($N=184$). All statistics are computed in the projected
(metre) frame; for fields far from the central meridian the UTM scale factor
varies negligibly across a single field.

\paragraph{Local order.} We use the Clark--Evans nearest-neighbour index
\citep{clark1954}, $R=\bar d_{\mathrm{NN}}/\big(1/(2\sqrt{\rho})\big)$, where
$\bar d_{\mathrm{NN}}$ is the mean nearest-neighbour distance and $\rho$ the
number density (estimated on the convex hull); $R<1$ indicates clustering, $R=1$
randomness, and $R>1$ regularity, with a normal-approximation $z$-score.

\paragraph{Long-range order.} We compute the radially-averaged structure factor
\begin{equation}
S(\mathbf{k})=\frac1N\Big|\sum_{j=1}^{N}e^{-i\mathbf{k}\cdot\mathbf{r}_j}\Big|^2,
\qquad \mathbf{k}=2\pi\big(n_x/L_x,\,n_y/L_y\big),
\label{eq:sk}
\end{equation}
on the allowed wavevectors of the bounding box, and report its low-$k$ limit
$S(k\!\to\!0)$ (the mean over $k<0.4\,k_\rho$, with $k_\rho=2\pi\sqrt{\rho}$). We
also compute the number variance $\sigma^2(R)$ from counts in random interior
disks of radius $R$, and its log--log slope: a slope near $2$ is area-like
(Poisson or clustered), a slope near $1$ is perimeter-like (hyperuniform).

\paragraph{Nulls.} Each field is compared against (i) complete spatial
randomness (CSR) in the same window; (ii) a random-sequential-adsorption (RSA)
hard-core process with exclusion distance matched to the field's
nearest-neighbour scale, representing the steric / lateral-inhibition
expectation; and (iii) an inhomogeneous-Poisson null that reproduces the field's
coarse density variations but is otherwise random, isolating the contribution of
heterogeneity. A Gaussian-perturbed lattice provides a disordered-hyperuniform
reference.

\section{Results}

\subsection{Local order spans clustered to regular}

The three fields differ markedly in local order (Table~\ref{tab:summary};
Fig.~\ref{fig:fields}). The northwestern Calabrian margin is strongly clustered
($R=0.47$, $z=-20.9$), its points falling into several patches. Offshore Zannone
Island is essentially random ($R=1.08$, $z=+1.6$), a compact, near-isotropic
cloud. The Malta Plateau is regular ($R=1.26$, $z=+6.7$) but forms an elongated,
diagonally-trending band (aspect ratio $2.5$) --- regular spacing along a
lineament rather than an isotropic field.

\subsection{The local order is genuine}

Comparison against the nulls (Table~\ref{tab:nulls}) shows these local signals
are real, not artefacts of window shape or density. The Calabrian
nearest-neighbour index ($0.47$) lies below even its inhomogeneous-Poisson null
($0.80$): the field is more clustered than its density variation alone would
produce, i.e.\ genuinely clustered into chains and patches. The Malta Plateau
index ($1.26$) exceeds not only CSR ($1.08$) but also the hard-core RSA null
($1.15$): its spacing is more regular than a simple steric gas, consistent with
an actively spaced (lateral-inhibition) lineament. Offshore Zannone tracks CSR
($1.08$ vs.\ $1.14$), showing no resolvable local exclusion.

\subsection{No field is hyperuniform}

The long-range picture is uniform across all three fields and is the central
result (Fig.~\ref{fig:orderplane}, Fig.~\ref{fig:nv}). Every field has
$S(k\!\to\!0)>1$ --- $2.5$ (Zannone), $19.4$ (Calabrian) and $5.4$ (Malta) ---
and a number variance that grows with window \emph{area} (log--log slope
$2.8$--$3.3$), not perimeter. The structure factor rises above unity at small
$k$ in every case and never shows the $S(k\!\to\!0)\to0$ suppression that defines
hyperuniformity (Fig.~\ref{fig:orderplane}b). In the order plane of local versus
long-range order (Fig.~\ref{fig:orderplane}a) the fields scatter widely along the
nearest-neighbour axis but all sit high on the fluctuation axis; the hyperuniform
region is empty.

Crucially, each field's long-range fluctuation level matches its own
inhomogeneous-Poisson null: $S(k\!\to\!0)=2.5$ versus $2.1$ (Zannone), $19.4$
versus $9.8$ (Calabrian), $5.4$ versus $6.7$ (Malta). The excess density
fluctuation is therefore largely the signature of density heterogeneity ---
patchiness, the elongated lineament, regional gradients --- and not of any
intrinsic long-range correlation. The only suppression of $S(k\!\to\!0)$ below
unity anywhere in the analysis appears in the hard-core RSA \emph{null}
($0.7$--$0.9$), never in the data.

\begin{table}[t]
\centering
\caption{The three analysed pockmark fields. $R$ is the Clark--Evans
nearest-neighbour index ($<1$ clustered, $>1$ regular); $d_0$ the median
nearest-neighbour distance; $S(k\!\to\!0)$ the low-$k$ structure factor; the last
column the log--log slope of the number variance ($\approx2$ area-like,
$\approx1$ hyperuniform).}
\label{tab:summary}
\begin{tabular}{lrrrrrrr}
\toprule
Field & $N$ & extent & aspect & $R$ & $z$ & $d_0$ & $S(k\!\to\!0)$ / NV slope\\
 & & (km) & & & & (m) & \\
\midrule
Offshore Zannone Island      & $112$ & $1.4$  & $1.8$ & $1.08$ & $+1.6$  & $42$  & $2.5$ / $2.8$\\
NW Calabrian margin          & $418$ & $37.5$ & $1.7$ & $0.47$ & $-20.9$ & $211$ & $19.4$ / $2.9$\\
Malta Plateau                & $184$ & $11.2$ & $2.5$ & $1.26$ & $+6.7$  & $103$ & $5.4$ / $3.3$\\
\bottomrule
\end{tabular}
\end{table}

\begin{table}[t]
\centering
\caption{Each field against its nulls: nearest-neighbour index $R$, low-$k$
structure factor $S(k\!\to\!0)$, and number-variance slope. Local order is real
(data $R$ departs from CSR in the expected direction and, for the clustered and
regular fields, beyond the heterogeneity and hard-core nulls), but the long-range
fluctuations of every field track its inhomogeneous-Poisson null --- none is
hyperuniform.}
\label{tab:nulls}
\begin{tabular}{llrrr}
\toprule
Field & Model & $R$ & $S(k\!\to\!0)$ & NV slope\\
\midrule
\multirow{4}{*}{Offshore Zannone}
 & data                  & $1.08$ & $2.54$ & $2.81$\\
 & CSR                   & $1.14$ & $1.00$ & $1.84$\\
 & RSA hard-core         & $1.31$ & $0.72$ & $1.84$\\
 & inhomogeneous-Poisson & $1.04$ & $2.05$ & $2.60$\\
\midrule
\multirow{4}{*}{NW Calabrian}
 & data                  & $0.47$ & $19.4$ & $2.90$\\
 & CSR                   & $1.03$ & $0.99$ & $1.87$\\
 & RSA hard-core         & $1.07$ & $1.02$ & $1.93$\\
 & inhomogeneous-Poisson & $0.80$ & $9.77$ & $3.32$\\
\midrule
\multirow{4}{*}{Malta Plateau}
 & data                  & $1.26$ & $5.36$ & $3.28$\\
 & CSR                   & $1.08$ & $1.04$ & $1.94$\\
 & RSA hard-core         & $1.15$ & $0.91$ & $1.93$\\
 & inhomogeneous-Poisson & $0.94$ & $6.73$ & $3.20$\\
\bottomrule
\end{tabular}
\end{table}

\begin{figure}[htbp]
\centering
\includegraphics[width=\textwidth]{fig_fields.pdf}
\caption{The three analysed pockmark fields in their projected (metre) frame,
centred. Local order ranges from clustered (Calabrian, multi-patch) through
near-random (Zannone) to regular but anisotropic (Malta Plateau, an elongated
lineament).}
\label{fig:fields}
\end{figure}

\begin{figure}[htbp]
\centering
\includegraphics[width=\textwidth]{fig_orderplane.pdf}
\caption{\textbf{(a)} Order plane: nearest-neighbour index (local order) versus
$S(k\!\to\!0)$ (long-range order, log scale). The fields scatter along the local
axis but all lie well above the Poisson line on the fluctuation axis; the
hyperuniform region (bottom, $S(k\!\to\!0)\to0$) is empty. Reference markers:
CSR, a hard-core liquid (RSA), and a disordered-hyperuniform pattern.
\textbf{(b)} Radial structure factor: every field rises above $S=1$ at small $k$
(density heterogeneity) and none shows the $S(k\!\to\!0)\to0$ dip of
hyperuniformity. $k_\rho=2\pi\sqrt{\rho}$.}
\label{fig:orderplane}
\end{figure}

\begin{figure}[htbp]
\centering
\includegraphics[width=0.62\textwidth]{fig_numbervariance.pdf}
\caption{Number variance $\sigma^2(R)$ versus window radius. For every field
$\sigma^2(R)$ grows with window area (slope $\approx2$--$3$), parallel to the
Poisson reference, rather than with perimeter (slope $1$, the hyperuniform
reference). No field shows hyperuniform suppression of density fluctuations.}
\label{fig:nv}
\end{figure}

\FloatBarrier
\section{Discussion}

Across three Mediterranean pockmark fields spanning the full range of local
order, the long-range verdict is unanimous: none is hyperuniform. Local
regularity, where present, is genuine --- the Malta Plateau is more regularly
spaced than a hard-core gas, and the Calabrian margin is genuinely clustered ---
but it does not propagate to suppressed long-wavelength density fluctuations.
Instead, the large-scale fluctuations of every field are super-Poissonian and
are accounted for by density heterogeneity, as shown by their agreement with the
inhomogeneous-Poisson nulls. This is the same local-versus-long-range structure
seen in other self-organised geomorphic patterns: order that is real at the scale
of neighbouring features but is overwritten at long range by the heterogeneity of
the natural substrate.

The single regular field sharpens the point. The Malta Plateau's high
nearest-neighbour index might, on a local statistic alone, be read as evidence of
an ordered, possibly hyperuniform, self-organised field. The structure factor and
number variance show otherwise: the regularity is one-dimensional, expressed
along an elongated lineament (aspect ratio $2.5$), and the isotropically-averaged
long-range fluctuations are if anything the largest of the three fields. Regular
nearest-neighbour spacing of the kind attributed to lateral inhibition
\citep{hammer2009} can thus coexist with strongly non-hyperuniform long-range
behaviour; the two must be measured separately.

We are explicit about what is and is not new. That pockmark fields can be
clustered along faults or regularly spaced by lateral inhibition is established
\citep{hammer2009,spatola2025}, as is the regional, fault-related distribution of
the Mediterranean database used here. The contribution of this note is the
hyperuniformity-class characterisation --- the structure factor, number variance
and matched nulls --- and the resulting verdict that the order in these fields is
local and heterogeneity-limited, not hyperuniform. Several caveats bound the
result. The database is a literature compilation of varying per-field
completeness; we mitigated this by selecting contiguous, artefact-free fields,
but completeness cannot be guaranteed, and incompleteness would, if anything, add
spurious fluctuations rather than suppress them, so it cannot mask a true
hyperuniform signal. The field sizes ($N\sim10^2$) limit the lowest accessible
wavenumber, so we claim only the absence of hyperuniformity, not a precise
fluctuation exponent. Finally, the analysis is isotropically averaged; the Malta
Plateau shows that directional analysis is the natural next step for lineated
fields.

\section{Conclusion}

Submarine pockmark fields exhibit local spatial order ranging from clustered to
regular, but a structure-factor and number-variance analysis of representative
Mediterranean fields finds no hyperuniformity: long-range density fluctuations
are super-Poissonian and governed by substrate heterogeneity, and the one
genuinely regular field is an anisotropic lineament. Nearest-neighbour
regularity in a fluid-escape field should therefore not be taken as evidence of
long-range order without a structure-factor test.

\section*{Data and code availability}
The pockmark geodatabase is from \citet{spatola2025}, Zenodo
\href{https://doi.org/10.5281/zenodo.15425907}{10.5281/zenodo.15425907}
(CC-BY-4.0). Analysis code (a pure-Python shapefile reader, the point-pattern
diagnostics, and the field-selection and figure scripts), the derived per-field
coordinates and all result tables are available at
\url{https://github.com/Ruqing1963/pockmark-field-order} and archived on Zenodo
(DOI assigned on release).

\section*{Acknowledgements}
This work uses the open pockmark geodatabase of Spatola and colleagues; their
compilation made this analysis possible.

\begin{thebibliography}{99}

\bibitem[Clark \& Evans(1954)]{clark1954}
P.~J. Clark and F.~C. Evans.
Distance to nearest neighbor as a measure of spatial relationships in
populations. \textit{Ecology} \textbf{35}, 445--453 (1954).

\bibitem[Hammer(2009)]{hammer2009}
\O.~Hammer.
Pattern formation by local amplification and lateral inhibition: examples from
biology and geology.
\textit{Eur. Phys. J. Special Topics} \textbf{178}, 5--12 (2009).

\bibitem[Hovland \& Judd(1988)]{hovland1988}
M.~Hovland and A.~G. Judd.
\textit{Seabed Pockmarks and Seepages: Impact on Geology, Biology and the Marine
Environment}. Graham \& Trotman, London (1988).

\bibitem[Jiao et al.(2014)]{jiao2014}
Y.~Jiao, T.~Lau, H.~Hatzikirou, M.~Meyer-Hermann, J.~C. Corbo, and S.~Torquato.
Avian photoreceptor patterns represent a disordered hyperuniform solution to a
multiscale packing problem. \textit{Phys. Rev. E} \textbf{89}, 022721 (2014).

\bibitem[King \& MacLean(1970)]{king1970}
L.~H. King and B.~MacLean.
Pockmarks on the Scotian Shelf.
\textit{Geol. Soc. Am. Bull.} \textbf{81}, 3141--3148 (1970).

\bibitem[Spatola et al.(2025)]{spatola2025}
D.~Spatola, M.~Rovere, D.~Casalbore, and F.~L. Chiocci.
Pockmarks of the Mediterranean region seas: a comprehensive geodatabase for
marine geomorphological analysis. \textit{Sci. Data} \textbf{12}, 1049 (2025).

\bibitem[Torquato(2018)]{torquato2018}
S.~Torquato.
Hyperuniform states of matter.
\textit{Phys. Rep.} \textbf{745}, 1--95 (2018).

\bibitem[Torquato \& Stillinger(2003)]{torquato2003}
S.~Torquato and F.~H. Stillinger.
Local density fluctuations, hyperuniformity, and order metrics.
\textit{Phys. Rev. E} \textbf{68}, 041113 (2003).

\end{thebibliography}

\end{document}
